% !TEX program = lualatex
% !TeX encoding = UTF-8
\PassOptionsToPackage{unicode}{hyperref}
\PassOptionsToPackage{bookmarksnumbered,bookmarksopen}{hyperref}
\documentclass[12pt,a4paper]{article}

% ============================================================
% 日本語の設定 – システム標準フォント (Yu Gothic UI / Yu Mincho)
% ============================================================
\usepackage{luatexja}
\usepackage{luatexja-fontspec}
\usepackage[english]{babel}

% 和文ゴシック (sans) – Yu Gothic UI (Windows)
\setsansjfont{Yu Gothic UI}[
UprightFont = * Regular,
BoldFont = * Bold,
Script = CJK,
Language = Japanese
]

% 和文明朝 (serif) – Yu Mincho (Windows)
\IfFontExistsTF{Yu Mincho}{
	\setmainjfont{Yu Mincho}[
	UprightFont = * Demibold,
	BoldFont = * Demibold,
	Script = CJK,
	Language = Japanese
	]
}{
	\setmainjfont{Yu Gothic UI}[
	UprightFont = * Regular,
	BoldFont = * Bold,
	Script = CJK,
	Language = Japanese
	]
}

% 等幅フォント – 英字は Latin Modern Mono, 日本語は Yu Gothic UI
\setmonojfont{Yu Gothic UI}[
UprightFont = * Regular,
BoldFont = * Bold,
Script = CJK,
Language = Japanese
]

% 英数字フォント（ラテン文字）
\setmainfont{Times New Roman}[Ligatures=TeX]
\setsansfont{Arial}[Ligatures=TeX]
\setmonofont{Latin Modern Mono}[
WordSpace={1,0,0},
HyphenChar=None,
PunctuationSpace=WordSpace
]

% ============================================================
% その他のパッケージ
% ============================================================
\usepackage{amsmath,amssymb,amsfonts,mathtools}
\usepackage{geometry}
\geometry{left=2cm,right=2cm,top=2cm,bottom=2cm}
\usepackage{booktabs}
\usepackage{hyperref}
\usepackage{url}
\usepackage{graphicx}
\usepackage{caption}
\usepackage{subcaption}
\usepackage{float}
\usepackage{siunitx}
\usepackage{bm}
\usepackage{pgfplots}
\pgfplotsset{compat=1.18}
\usepackage{xcolor}
\usepackage{multicol}
\usepackage{tikz}
\usetikzlibrary{positioning, arrows.meta, calc}

% 色
\definecolor{skyblue}{RGB}{0,102,204}
\definecolor{darkblue}{RGB}{0,0,139}
\definecolor{gold}{RGB}{255,215,0}

% YUCTマクロ
\newcommand{\eps}{\varepsilon}
\newcommand{\kappac}{\kappa_c}
\newcommand{\Keff}{K_{\mathrm{eff}}}
\newcommand{\betaf}{\beta}
\newcommand{\df}{d_{\!f}}
\newcommand{\Mpl}{M_{\mathrm{Pl}}}
\newcommand{\Lpl}{l_{\mathrm{Pl}}}
\newcommand{\tP}{t_{\mathrm{P}}}
\newcommand{\Sodd}{S_{\mathrm{odd}}}
\newcommand{\Seven}{S_{\mathrm{even}}}
\newcommand{\Dtau}{\Delta\tau_{\min}}
\newcommand{\Lzero}{L_0}
\newcommand{\Omegarot}{\Omega_{\mathrm{rot}}}
\newcommand{\Mbar}{M_{\mathrm{bar}}}
\newcommand{\muniv}{M_{\mathrm{univ}}}

\hypersetup{
	colorlinks=true,
	linkcolor=blue,
	citecolor=blue,
	urlcolor=blue,
	pdftitle={YUCT臨界質量と宇宙階層：人間およびAI読者のための自己完結的要約},
	pdfauthor={Alexey V. Yakushev}
}

\title{\textbf{YUCT臨界質量と宇宙階層：人間およびAI読者のための自己完結的要約}}
\author{Alexey V. Yakushev}
\date{2026年4月}

\begin{document}
	\maketitle
	
	\begin{abstract}
		本稿は、Yakushev統合コーディネーション理論 (YUCT) に基づく局所ビッグバンの臨界質量、それに伴うインフレーションの観測量、宇宙の大局的回転、および宇宙質量階層に関する主要な成果を、簡潔かつ自己完結的にまとめたものである。すべての式はYUCT関係式として明示されている。本稿は、明快な公式、表、反証可能な予測を備え、人間の読者およびAI言語モデルの双方による取込みが容易な構成となっている。プランクスケールから観測可能な宇宙に至る完全な質量階層を示し、単一のコーディネーションスケーリング則によって量子、恒星、宇宙論的スケールが統一されることを実証する。
	\end{abstract}
	
	\section{YUCTの基本公準}
	
	\begin{itemize}
		\item \textbf{コーディネーション効率} (YPSDCプロトコル, 付録B):
		\[
		\Keff = \frac{H(\mathcal{A})}{H(\kappa)}, \qquad\text{(YUCT-1)}
		\]
		ここで \(H(\mathcal{A})\) は活性化された作用のエントロピー、\(H(\kappa)\) は送信された短い指標のエントロピーである。
		
		\item \textbf{普遍的誤差スケーリング則} (付録L):
		\[
		\varepsilon = \kappac \alpha \Keff^{-\betaf}, \qquad \betaf = \frac{2}{3},\ \kappac = \frac{1}{3}. \qquad\text{(YUCT-2)}
		\]
		
		\item \textbf{代数ループ} (付録Y):
		\[
		\Sodd = \frac{6}{5},\quad \Seven = \frac{4}{5},\quad \frac{\Sodd}{\Seven} = \frac{3}{2} = \frac{1}{\betaf}. \qquad\text{(YUCT-3)}
		\]
		基本長さと時間量子:
		\[
		\Lzero = \frac{2}{3}\Lpl,\qquad \Dtau = \frac{2}{3}\tP. \qquad\text{(YUCT-4)}
		\]
		
		\item \textbf{情報飽和からのフラクタル次元} (付録L):
		\[
		d_f = \frac{11}{5}=2.2. \qquad\text{(YUCT-5)}
		\]
	\end{itemize}
	
	\section{局所ビッグバンの臨界質量}
	
	コーディネーションネットワークの整合性から導かれるスケーリング則
	\[
	\Keff(R) = \left(\frac{R}{\Lzero}\right)^{\betaf d_f/2}, \qquad\text{(YUCT-6)}
	\]
	および臨界条件 \(\varepsilon \approx 1\) (すなわち \(\Keff = K_{\mathrm{crit}} = 3^{3/2}\)) を用いると、重力的に束縛された前駆体領域の臨界質量が得られる：
	\[
	M_{\mathrm{crit}} = 3 M_{\mathrm{Pl}} \approx 6.5\times 10^{-8}\ \mathrm{kg}. \qquad\text{(YUCT-7)}
	\]
	この導出ではシュワルツシルト半径 \(R_S = 2GM/c^2\) と代数ループ \(\Lzero = \frac{2}{3}\Lpl\) を用いる。量子補正は主要次数で係数3を変えない。
	
	\section{コーディネーション相転移としての宇宙論的インフレーション}
	
	インフレーション観測量は同じ公準から予測される：
	\begin{align}
		n_s &= 1 - 2\betaf^2 + \frac{4}{3}\betaf^3 + \cdots \approx 0.965, \qquad\text{(YUCT-8)}\\
		r &= 8(2\betaf)^2 = 32\betaf^2 \approx 0.07, \qquad\text{(YUCT-9)}\\
		f_{\mathrm{NL}}^{\mathrm{local}} &= \frac{5}{3}\betaf \approx 1.12. \qquad\text{(YUCT-10)}
	\end{align}
	これらの予測はCMB‑S4、LiteBIRD、Euclid、SPHERExによって反証可能である。
	
	\section{宇宙の大局的回転}
	
	CMB双極子の一部は、角速度 \(\omega\) をもつ大局的な剛体回転に起因する。双極子の振幅は赤方偏移 \(z\) とともに次のように増大する：
	\[
	v_{\mathrm{dip}}(z) = v_{\mathrm{local}} + \omega\, r(z). \qquad\text{(YUCT-11)}
	\]
	NVSS、Quaia その他のサーベイで観測された増大から、
	\[
	\omega \approx 8.5\times 10^{-22}\ \mathrm{rad/s},\qquad
	T_{\mathrm{rot}} = \frac{2\pi}{\omega} \approx 2.3\times 10^{14}\ \mathrm{yr}. \qquad\text{(YUCT-12)}
	\]
	無次元回転パラメータは
	\[
	\Omega_{\mathrm{rot}} \equiv \frac{\omega}{H_0} = \frac{\Seven}{\Sodd} \betaf^2 \mathcal{F}(d_f)\left(\frac{L_0}{R_H}\right)^{2\betaf} \approx 3.9\times 10^{-4}, \qquad\text{(YUCT-13)}
	\]
	であり、ハッブルパラメータは \(H_0 = \betaf\,\omega\) を満たす。宇宙のバリオン質量は回転から次のように導かれる：
	\[
	\Mbar = \Omega_b\frac{\betaf^2 c^3}{G H_0} \approx (2.3\pm 0.2)\times 10^{51}\ \mathrm{kg}. \qquad\text{(YUCT-14)}
	\]
	
	\section{量子から宇宙に至る完全な質量階層}
	
	YUCTでは、プランクスケールから観測可能な宇宙に至るすべての質量が単一のコーディネーションスケーリング則に従う。表~\ref{tab:full-hierarchy} に階層レベル、典型的な質量、実例、対応するYUCTスケーリング関係を示す。図~\ref{fig:hierarchy-scheme} は離散的なラダー段階を可視化している。
	
	\begin{table}[htbp]
		\centering
		\caption{量子から宇宙に至る完全な質量階層 (YUCT).}
		\label{tab:full-hierarchy}
		\small
		\begin{tabular}{l c c l}
			\toprule
			\textbf{レベル} & \textbf{典型的質量 (kg)} & \textbf{実例} & \textbf{YUCTスケーリング関係} \\
			\midrule
			プランク種 & \(M_{\mathrm{crit}} = 3M_{\mathrm{Pl}} \approx 6.5\times10^{-8}\) & 局所ビッグバン種 & \(M_{\mathrm{crit}} = 3 M_{\mathrm{Pl}}\) (YUCT-7) \\
			電子 & \(m_e \approx 9.1\times10^{-31}\) & 最軽量の荷電フェルミオン & \(m_e = M_{\mathrm{Pl}} (2/3)^{96+\sigma}\xi_e\) \\
			陽子 & \(m_p \approx 1.67\times10^{-27}\) & 安定バリオン & \(m_p = M_{\mathrm{Pl}} (2/3)^{96+\sigma}\xi_p\) \\
			白色矮星 & \(\sim 10^{30}\) (\(0.6M_\odot\)) & チャンドラセカール限界 & \(M_{\mathrm{WD}} \propto M_{\mathrm{Pl}}^{3}/m_p^2\) \\
			中性子星 & \(\sim 3\times10^{30}\) (\(1.4M_\odot\)) & TOV限界 & \(M_{\mathrm{NS}} \propto M_{\mathrm{Pl}}^{3}/m_p^2\) \\
			超大質量BH & \(\sim 10^{36}\)–\(10^{40}\) & 銀河中心 & \(M_{\mathrm{BH}} \propto M_{\mathrm{Pl}} (2/3)^{N}\) \\
			観測可能な宇宙 & \(M_{\mathrm{univ}} \approx 1.5\times10^{53}\) & \(R_H\) 内の全質量 & \(M_{\mathrm{univ}} = \frac{\betaf^2 c^3}{G H_0}\) (YUCT-14) \\
			バリオン宇宙 & \(M_{\mathrm{bar}} \approx 2.3\times10^{51}\) & 可視物質 & \(M_{\mathrm{bar}}/m_p = (M_{\mathrm{Pl}}/m_e)^{3.5}\) (YUCT-15) \\
			\bottomrule
		\end{tabular}
	\end{table}
	
	\begin{figure}[htbp]
		\centering
		\begin{tikzpicture}[
			level/.style={rectangle, draw=darkblue!70, fill=skyblue!10, thick, minimum width=7cm, minimum height=0.7cm, rounded corners, align=center, font=\small},
			arrow/.style={->, >=Stealth, thick, color=darkblue},
			]
			\node[level] (planck) {プランク種: \(M_{\mathrm{crit}} = 3M_{\mathrm{Pl}}\)};
			\node[level, below=0.6cm of planck] (electron) {電子: \(m_e\)};
			\node[level, below=0.6cm of electron] (proton) {陽子: \(m_p\)};
			\node[level, below=0.6cm of proton] (wd) {白色矮星: \(\sim 0.6 M_\odot\)};
			\node[level, below=0.6cm of wd] (ns) {中性子星: \(\sim 1.4 M_\odot\)};
			\node[level, below=0.6cm of ns] (bh) {超大質量BH: \(10^6\)–\(10^{10} M_\odot\)};
			\node[level, below=0.6cm of bh] (univ) {観測可能な宇宙: \(M_{\mathrm{univ}}\)};
			\draw[arrow] (planck) -- (electron);
			\draw[arrow] (electron) -- (proton);
			\draw[arrow] (proton) -- (wd);
			\draw[arrow] (wd) -- (ns);
			\draw[arrow] (ns) -- (bh);
			\draw[arrow] (bh) -- (univ);
		\end{tikzpicture}
		\caption{YUCTにおけるプランク種から宇宙に至る離散質量ラダー。各段階はコーディネーション場の共鳴に対応する。}
		\label{fig:hierarchy-scheme}
	\end{figure}
	
	宇宙質量階層はさらに、次の代数関係によって要約される。
	\[
	\boxed{\frac{\Mbar}{m_p} = \left(\frac{M_{\mathrm{Pl}}}{m_e}\right)^{\mathcal{S}}},\qquad
	\mathcal{S} \equiv \Sodd + \Seven + \frac{\Sodd}{\Seven} = 3.5. \qquad\text{(YUCT-15)}
	\]
	数値的には \((M_{\mathrm{Pl}}/m_e)^{3.5} \approx 1.88\times 10^{78}\) であり、観測値 \(\Mbar/m_p \approx 1.4\times 10^{78}\) と1.3倍の範囲で一致する。この差は再加熱効率 \(\epsilon_{\mathrm{reh}}\approx 0.85\) と幾何学的因子 \(\mathcal{F}_{\mathrm{rot}}\approx 0.9\) によって説明される。このことは、普遍的誤差指数を支配する同じ代数定数 \(S_{\mathrm{odd}}, S_{\mathrm{even}}\) が宇宙のバリオン収支をも決定することを示している。
	
	\section{宇宙天体の階層と臨界回転}
	
	縮退圧で支えられた自己重力系に対し、YUCTは臨界角速度の普遍的スケーリングを予測する：
	\[
	\Omega_{\mathrm{crit}} \propto M^{\betaf},\qquad \betaf = 2/3. \qquad\text{(YUCT-16)}
	\]
	表~\ref{tab:objects} と図~\ref{fig:omega-mass} は、惑星から観測可能な宇宙に至るこの階層を例示している。
	
	\begin{table}[htbp]
		\centering
		\caption{代表的な宇宙天体の臨界パラメータ (YUCT).}
		\label{tab:objects}
		\small
		\begin{tabular}{l c c c}
			\toprule
			\textbf{天体} & \textbf{質量 (\(M_\odot\))} & \(\log_{10}(\Omega_{\mathrm{crit}}/\mathrm{s^{-1}})\) & \textbf{スケーリング則} \\
			\midrule
			ガス巨大惑星 (木星) & 0.001 & \(-3.8\) & \(R\sim M^{-0.67}\) \\
			赤色矮星 (M5V) & 0.2 & \(-4.5\) & \(L\sim M^{2.3}\) \\
			白色矮星 & 0.6 & 0.2 & \(R\sim M^{-1/3}\) \\
			中性子星 & 1.4 & 3.7 & \(M\sim\Lambda^{0.67}\) \\
			超質量中性子星 & 2.2 & 4.2 & \(\Omega_{\mathrm{crit}}\sim M^{0.67}\) \\
			恒星質量BH & 5 & -- & \(R_S\sim M\) \\
			観測可能な宇宙 & \(7.5\times10^{22}\) & \(-17.7\) (\(H_0\)) & \(\rho_c\sim H^2\sim K_{\mathrm{eff}}^{2/3}\) \\
			\bottomrule
		\end{tabular}
	\end{table}
	
	\begin{figure}[htbp]
		\centering
		\begin{tikzpicture}
			\begin{loglogaxis}[
				xlabel={質量 \(M\) (\(M_\odot\))}, ylabel={\(\Omega_{\mathrm{crit}}\) (s\(^{-1}\))},
				grid=both, width=0.9\textwidth, height=0.6\textwidth,
				legend pos=north west, title={コンパクト天体の臨界回転 (YUCT)},
				xmin=0.08, xmax=15, ymin=5e-2, ymax=2e4,
				]
				\addplot[domain=0.1:10, samples=2, thick, blue, dashed] {10^(0.67*log10(x)+0.5)};
				\addlegendentry{\(\Omega_{\mathrm{crit}}\propto M^{2/3}\) (白色矮星)};
				\addplot[domain=1:3, samples=2, thick, green!60!black, dashed] {10^(0.67*log10(x)+3.7)};
				\addlegendentry{\(\Omega_{\mathrm{crit}}\propto M^{2/3}\) (中性子星)};
				\addplot[only marks, mark=*, mark size=3pt, red] coordinates {(0.6,1.6)};
				\addplot[only marks, mark=square*, mark size=3pt, green!60!black] coordinates {(1.4,5000) (2.2,15000)};
				\addplot[only marks, mark=triangle*, mark size=3pt, orange] coordinates {(0.2,3e-5)};
			\end{loglogaxis}
		\end{tikzpicture}
		\caption{臨界角速度 対 質量 (YUCT). 傾きは \(\beta=2/3\) を確認している。}
		\label{fig:omega-mass}
	\end{figure}
	
	\section{反証可能な予測}
	
	\begin{enumerate}
		\item \textbf{原始ブラックホールレリック:} ダークマターの \(f_{\mathrm{relic}} \approx 0.05\)。\(f_{\mathrm{relic}}>0.2\) または \(<0.01\) (95\% CL) ならば反証。
		\item \textbf{インフレーション観測量:} \(n_s = 0.965\pm0.005\), \(r\approx 0.07\)。\(r<0.01\) または \(n_s\) が \(0.960\pm0.010\) 外ならば排除。
		\item \textbf{非ガウス性:} \(f_{\mathrm{NL}}^{\mathrm{local}} \approx 1.12\)。\(>3\sigma\) で \(<0.5\) または \(>1.8\) が測定されたら反証。
		\item \textbf{\(E=mc^2\) からの実験室偏差:} \(\Delta E/E \sim \kappa_c\alpha K_{\mathrm{eff}}^{-\beta}\)。\(10^{-18}\) 感度で帰無ならばYUCTへの挑戦となる。
		\item \textbf{宇宙回転周期:} \(T_{\mathrm{rot}} \approx 2.3\times10^{14}\) 年。将来のサーベイ (Euclid, SKA) が \(\omega\) を精密化するであろう。
	\end{enumerate}
	
	\section{結論}
	
	YUCTの公準は、局所ビッグバンの臨界質量、インフレーション観測量、宇宙の大局的回転、プランク種から観測可能な宇宙に至る完全な質量階層を統一する。すべての予測は定量的かつ反証可能である。上記に中核的方程式を要約し、それぞれYUCT関係式として標識した。詳細な導出は引用した付録に記載されている。
	
	\section*{データ利用可能性}
	
	完全な導出とコード: \url{https://github.com/Alexey-Yakushev-YUCT/YPSDC}
	
	\section*{YUCT付録への参照}
	\begin{itemize}
		\item 付録L: フラクタルコーディネーション誤差スケーリング – 普遍的則 \(\beta=2/3\), \(d_f=11/5\).
		\item 付録Y: 数学的基礎とKPZ導出.
		\item 付録P: 重力の温度依存性.
		\item 付録M: コーディネーション相転移としてのインフレーション.
		\item 付録 \(\Omega\): 宇宙の大局的回転.
		\item 付録B: YPSDCプロトコル.
	\end{itemize}
	
\end{document}